\begin{abstract}
  Non-invasive decoding of naturalistic stimuli has advanced rapidly over the past decade, with significant breakthroughs in fMRI and MEG across language, imagery, and speech domains. EEG, however, has largely been excluded from this trajectory. Naturalistic EEG decoding remains tethered to linear Temporal Receptive Field (TRF) models, while emerging self-supervised foundation models have focused almost exclusively on short-window clinical and BCI benchmarks. We introduce Cine-JEPA, a Joint-Embedding Predictive Architecture that learns EEG representations by predicting the latents of masked spatiotemporal blocks of neural activity during naturalistic movie-watching. Pretrained on the Healthy Brain Network corpus with a multi-block masking objective, the model learns to infer the latents of held-out blocks from the surrounding visible context rather than reconstructing the raw signal. Across a hierarchy of decoding probes—spanning low-level visual features (luminance, contrast), narrative events, and temporal position in the movie—Cine-JEPA substantially outperforms TRF baselines, fully supervised end-to-end models, and both linear-probe and full-finetuning evaluations of leading EEG foundation models. The representations additionally transfer to subject trait prediction, indicating gains beyond stimulus-locked decoding. Ablations over mask-block scale (short vs. long), context geometry (number and duration of temporal windows), and regularization isolate the contribution of masked latent prediction. Our results establish hierarchical decoding from naturalistic movie EEG as a discriminating evaluation regime for EEG self-supervised learning, and masked latent-block prediction as a competitive paradigm for neural representation learning more broadly.

\end{abstract}

\section{Introduction}
 
Naturalistic stimulus decoding---extracting semantic, perceptual, and temporal
structure from brain responses to continuous, ecologically valid inputs like
movies and natural speech---has rapidly transformed cognitive neuroscience over
the past decade. In fMRI, naturalistic paradigms have become the standard for
individual-difference and developmental studies \citep{vanderwal2019, finn2022},
and recent work has decoded continuous narrative language \citep{tang2023} and
natural images \citep{scotti2024} from cortical responses with striking fidelity.
MEG has followed a parallel trajectory, with dedicated benchmarks for
naturalistic speech \citep{gwilliams2023, ozdogan2025} and contrastive
self-supervised methods that decode perceived speech directly from non-invasive
recordings \citep{defossez2023}. Intracranial work has likewise been formalized
into multitask naturalistic decoding benchmarks \citep{neuroprobe2025}. EEG, by
contrast, has been largely left out of this trajectory---despite being the most
accessible non-invasive modality and the only one with population-scale
developmental movie-watching corpora \citep{alexander2017}.
 
Two factors have produced this asymmetry. First, naturalistic EEG decoding has
remained methodologically anchored to linear Temporal Receptive Field (TRF)
models from the early 2010s \citep{lalor2010, crosse2016, diliberto2018}. TRFs
are well-understood and interpretable but represent a fundamentally linear,
fully supervised approach to a problem the rest of the brain-decoding field has
moved past. Second, the self-supervised learning revolution that has produced
large EEG foundation models---BENDR \citep{kostas2021}, BIOT \citep{yang2023},
LaBraM \citep{jiang2024}, and the recent EEG-VJEPA \citep{hojjati2025}---has
been channeled almost exclusively into short-window clinical and BCI benchmarks:
sleep staging, abnormality detection, motor imagery, and seizure detection,
a pattern that fits the broader ``static encoder'' tendency in biosignal
foundation models recently characterized by \citet{wu2026}. The pretext tasks
(contrastive learning over augmented windows, masked reconstruction over
symmetric tubelets) and the evaluation regimes (categorical labels averaged
across trials) co-evolved with this paradigm, leaving the question of how
self-supervised learning should be designed for continuous, stimulus-driven
EEG essentially unaddressed.
 
We introduce \textbf{Cine-JEPA}, a Joint-Embedding Predictive Architecture
\citep{lecun2022, assran2023, bardes2024} that targets naturalistic-stimulus EEG
directly. Pretrained on the Healthy Brain Network corpus, the model uses a
multi-block masking objective: a context encoder embeds visible spatiotemporal
blocks of EEG, a target encoder updated by exponential moving average embeds
held-out blocks, and a predictor maps from context embeddings to target
embeddings. Unlike masked-reconstruction pretraining, no signal-space decoding
is required; gradient flow is shaped by latent prediction loss alone. We
evaluate transfer to (i)~naturalistic movie decoding, where stimulus-locked
representations must generalize across subjects and clips, and (ii)~subject
trait prediction, where the same representations must support
individual-difference inference. On both, Cine-JEPA substantially outperforms
TRF baselines, fully supervised end-to-end models, and both linear-probe and
full-finetuning evaluations of leading EEG foundation models.
 
Our contributions are as follows. \textbf{(1)}~We establish naturalistic movie
decoding from EEG as an evaluation regime that discriminates between
self-supervised methods in ways invisible under conventional clinical and BCI
benchmarks. \textbf{(2)}~We propose Cine-JEPA, a multi-block masked
latent-prediction architecture for continuous EEG, and show that its
representations outperform TRF baselines, supervised end-to-end models, and
linear-probe and full-finetuning evaluations of BENDR, BIOT, LaBraM, and
EEG-VJEPA on movie decoding. \textbf{(3)}~We demonstrate that Cine-JEPA
representations transfer beyond stimulus-locked decoding to subject trait
prediction, indicating that masked latent prediction captures structure
relevant to individual differences in addition to time-locked stimulus content.
\textbf{(4)}~Through configuration sweeps over mask-block scale, context
geometry (number and duration of visible windows), and regularization, we
identify the design settings under which masked latent-block prediction
succeeds on continuous EEG, and characterize how these effective configurations
depart from the defaults inherited from image- and video-domain JEPAs.